% Root post-CP22 candidate. Complete direct proof, independent audits pending.
\subsection{A direct exclusion of three total trigonal ramifications}

\begin{lemma}[The trigonal pencil on a nonordinary genus-two curve]
\label{lem:root-nonordinary-trigonal-three-excluded-candidate}
\status{Candidate; complete direct proof and universal identities written, independent audits pending.}
Let $C$ be a smooth projective genus-two curve over an algebraically
closed field $k$ of characteristic two, with two-rank zero or one.
No degree-three morphism $C\to\mathbf P^1_k$ has three distinct
totally ramified points.
\end{lemma}
\begin{proof}
Suppose such a morphism exists, and choose one totally ramified point
$Q$. It cannot be Weierstrass: the gap sequence at a genus-two
Weierstrass point is $1,3$, so no rational function has its sole pole
of exact order three there. Choose a Weierstrass point $W$ over
infinity for the hyperelliptic coordinate. The already proved
nonordinary normal forms allow the equation to have $h$ constant
or linear and a degree-five right side. For two-rank one we choose
the Weierstrass point with the order-three Artin--Schreier pole over
infinity, as in the preceding normal-form proof. Since $Q$ is not
Weierstrass, it is a finite point with $h(Q)\ne0$.

Translate its hyperelliptic coordinate to zero, translate its
$y$-coordinate to zero, and scale $y$ so that $h(0)=1$.
An additional change $y\mapsto y+v x^2$, with $v$ chosen to cancel
the degree-four coefficient, is possible over $k$ and does not move
$Q$. Thus we obtain
\[
 y^2+(1+t x)y=u x^5+a x^3+b x^2+c x,
 \qquad Q=(0,0),\quad u\ne0.                         \tag{RTG1}
\]
Here $t=0$ in two-rank zero and $t\ne0$ in two-rank one. These
changes are invertible on the function field and preserve the
smooth projective curve and its specified ramification question.
Put $e=b+c^2+ct$ and define
\[
                 z=\frac{y+1+(c+t)x+e x^2}{x^3}.       \tag{RTG2}
\]
At $Q$ the numerator is a unit, so $z$ has a pole of order three.
At its hyperelliptic conjugate $Q'=(0,1)$, the coefficient expansion
of (RTG1) gives
\[
        y=1+(c+t)x+(b+c^2+ct)x^2+O(x^3),
\]
so there is no pole there. There are no other finite possible poles.
At $W$, the pole orders of $x,y$ are two and five; hence $z$ has
a zero of order exactly one. Thus $z$ defines a separable degree-three
map with pole divisor $3Q$, and is unramified at $W$.
Riemann--Roch gives $h^0(3Q)=2$, so every degree-three map totally
ramified at $Q$ differs from $z$ by a target fractional linear
transformation. After moving that ramified value to infinity, the
change is affine. We may therefore count all the other totally
ramified fibres using $z$.

Substituting (RTG2) into (RTG1) and cancelling $x^3$ gives
\[
 z^2x^3+u x^2+(e^2+t z)x+z+a+te=0.                    \tag{RTG3}
\]
The value $z=0$ is not totally ramified, since its fibre contains
the unramified point $W$. If a nonzero finite value $z_0$ is totally
ramified, the cubic in $x$ in (RTG3) must be
$z_0^2(x-r)^3$. Its quadratic coefficient forces $r=u/z_0^2$.
The remaining coefficients give
\[
 \begin{split}
 P(z_0)&=0,&P(z)&=t z^3+e^2z^2+u^2,\\
 Q_5(z_0)&=0,&Q_5(z)&=z^5+(a+te)z^4+u^3.
 \end{split}                                         \tag{RTG4}
\]
If $t=0$, then $P(z)=(ez+u)^2$ has at most one root. There
cannot be two further totally ramified values, proving the assertion
in this case.

Suppose $t\ne0$. Put
\[
 \begin{split}
 D&=e^4+ua+ute,& A&=te^2+u,\\
 R(z)&=t^2z^2+uz+D,&
 R_1&=t^2D+uA,\qquad R_0=u^2t^3+AD,\\
 R_2&=A^2+ut^5.
 \end{split}                                         \tag{RTG5}
\]
The identity $P(z)^2+uQ_5(z)=z^4R(z)$ shows that every
nonzero common root in (RTG4) is a root of $R$.
Two further totally ramified points would give two distinct such
values. The quadratic $R$ is separable, since its derivative is
$u\ne0$, so it would divide $P$. The exact polynomial division is
\[
              t^3P(z)+(t^2z+A)R(z)=R_1z+R_0.
                                                               \tag{RTG6}
\]
Consequently $R_1=R_0=0$. The identity
$t^2R_0+AR_1=uR_2$ then gives $R_2=0$.

These conditions contradict smoothness of (RTG1). Indeed define
\[
                  S=u^2+a^2t^4+ut^5+at^7+et^8.
\]
There is the universal polynomial identity over $\mathbf F_2$
\[
                  u^2S=R_1^2+ut^5R_1+e^4t^2R_2.       \tag{RTG7}
\]
Thus $S=0$. If $f(x)$ denotes the right side of (RTG1), direct
substitution gives
\[
          S=t^8\bigl(f'(1/t)^2+t^2 f(1/t)\bigr).       \tag{RTG8}
\]
At $x=1/t$ the coefficient of $y$ is zero. The unique point with
$y^2=f(1/t)$ is singular precisely when
$ty+f'(1/t)=0$, which is equivalent to the vanishing in (RTG8).
This is a singular point of the affine model, contradicting its
smoothness. The transformations to (RTG1) were invertible, so this
is a contradiction to the original hypothesis.

The exact program
\path{code/nonordinary_trigonal_direct_certificate.py}
checks the elimination, the two triple-root conditions, both
remainder identities and the singularity identity as polynomial
identities over $\mathbf F_2$, including the separate $t=0$ case.
It does not replace the Riemann--Roch or ramification arguments.
\end{proof}

\begin{corollary}[No theta triples under a group killed by three]
\label{cor:root-nonordinary-theta-no-three-triples-candidate}
\status{Candidate consequence; independent audits pending.}
Let $(B,\Theta)$ be a nonproduct nonordinary principally polarized
abelian surface over $k$ as above, and let $K\subset B[3](k)$ be
any subgroup. No three distinct points of its smooth principal
theta curve have the same image in $B/K$. No isotropy assumption
on $K$ is needed for this assertion.
\end{corollary}
\begin{proof}
Identify $B$ with the Jacobian of its genus-two theta curve $C$,
using the already proved principal-surface dichotomy. Three such
points $P,Q,R$ would satisfy
$3P\sim3Q\sim3R$ as divisors on $C$. A rational function with
divisor $3P-3Q$ gives a degree-three map totally ramified at
$P,Q$. Since $h^0(3Q)=2$, the function with divisor $3R-3Q$
is an affine linear combination of that function and the constant
function. The same map is therefore totally ramified at $R$.
This contradicts the lemma.
\end{proof}

\begin{remark}[Scope and independent route]
\label{rem:root-trigonal-direct-scope-candidate}
\status{Candidate scope statement.}
This is a direct argument on the actual nonordinary characteristic-two
normal forms. It neither imports a generic complex nodality theorem
nor asserts a condition for higher odd kernel exponent. A separate
Galois-closure and Cartier argument has been proposed independently;
the direct proof above has no dependence on that argument or on its
ramification classification. The existence of full node permutation
groups and gluing translations remains a separate monodromy question.
\end{remark}
